% page 529 of the facsimile (image p-528 of the scan)
% block 165.1 x 237.7 mm ; left margin 36.9 mm
\pgc{15.240mm}{0.000mm}{
\l{\cel{52}521}
\l{}
\l{serpo. Larja lamo krecento-forma, di qua la aristo konkava}
\l{\cel{3}esas tranchiva, kun mancho kurta, uzata da la gardenisti,}
\l{\cel{3}hakisti, por emundar la arbori. - F.}
\l{}
\l{\sou{serpento.(zool.}) Reptero di qua la korpo esas tre oblonga, sen-}
\l{\cel{3}membra ; ula speci, venenoza ; qua su movas per la konkavaji}
\l{\cel{3}quin lu facas sur sulo. - L.\cel{1}\sou{serpens}. DEFIS.}
\l{}
\l{\sou{serpentino}. (\sou{mineral.}) Silikato di magnezio, kun makuli verda.}
\l{\cel{3}DEFIS.}
\l{}
\l{\sou{serpolo}. (\sou{bot.}) Planto labiea, aromatoza. - L. thymus serpillum.}
\l{\cel{3}EFIS.}
\l{}
\l{\sou{serumo}. (\sou{farmacio)}. Liquido aquatra, qua su separas de la parto}
\l{\cel{3}koagulinta di la sango di animali diversa, e preparita por}
\l{\cel{3}uzesor kom injektajo terapiala. - DEFIS.}
\l{}
\l{\sou{seruro}. (\sou{tekn.}) Aparato qua, esence, konsistas ye buxo fera}
\l{\cel{3}(\textquotedbl{}palastro\textquotedbl{}) quan onu aplikas ad-an pordo, e kovrilo di kofro,}
\l{\cel{3}e c.- peco movebla interne di ta buxo (\textquotedbl{}lango\textquotedbl{}) e peco kava}
\l{\cel{3}(\textquotedbl{}boko\textquotedbl{}) qua recevas la lango e fixigesis ube la klapo di la}
\l{\cel{3}pordo o di la kofro retrovenas por klozo ; fine, klefo, por}
\l{\cel{3}apertar o klozar la pordo, la kovrilo, e c. per igar la lango}
\l{\cel{3}enirar od ekirar la boko. - FIS.}
\l{}
\l{\sou{servar}. (\sou{trans.}) I. Igar su utila ad (ulu) po salario-quanto.}
\l{\cel{3}- II. (+\sou{serviar}) Prokurar (ad ulu) avantajo per sua interveno}
\l{\cel{3}personala. - DEFIS.}
\l{}
\l{\sou{servalo}. (\sou{zool.}) Tigro-kato di Afrika. -L.\cel{1}\sou{felis serva}). DEFRS.}
\l{}
\l{\sou{servico}. Ensemblo de la vazi (pladi, pladegi, glasi, tasi, boli,}
\l{\cel{3}e c.) e manjili (forci, kulieri, kulteli, e c.) o di la}
\l{\cel{3}linjo-tuki necesa por la repasto-tablo. - DEFIRS.}
\l{}
\l{\sou{servitudo}. (\sou{yurocienco}). Obligeso, debo, a qua submisesas ter-}
\l{\cel{3}proprietajo koncerne altra terproprietajo. - F.}
\l{}
\l{\sou{sesgar}. (\sou{trans}.)Entamar per deprenar parto di labordo. - S.}
\l{}
\l{\sou{sesiono.}\cel{2}Parto di la yaro, dum qua +seancas kontinue deliber-}
\l{\cel{3}antaro, tribunalo nepermananta. - DEFIRS.}
\l{}
\l{\sou{setono}. (\sou{medic.}) Mecho de kotono quan onu pasigas tra la pelo}
\l{\cel{3}e la celularo por durigar la funciono di la moyeno per qua}
\l{\cel{3}onu ekfluigas ulo mala e liberigas su de lu. - EFI.}
\l{}
\l{\sou{severa}. (\sou{ad, pri)}\cel{2}I. Sen indulgo pri la kulpi. - II. Qua obe-}
\l{\cel{3}dias strikte la regulo. - EFIS.}
\l{}
\l{\sou{seviciar}. (\sou{ad)}. (\sou{netrans}.)(\sou{yurocienco}). Agar violentoze : spozo,}
\l{\cel{3}a sua spozo ; patro, a sua filio ; mastro, a sua servisto.}
}
